\scanpage{91}
Ainsi, on obtient via le pol. car. modifié un hom
\[
(V,f)\longmapsto\frac{1}{\det(\mathrm{id}_V-ft)}=L(\unsure{V|N},f)
\]
\[
R(\deleted{\unsure{\mathbb{Z}}}\mathbb{N};k)\longrightarrow 1+k[[t]]^+
\]
\margin{en fait, à valeurs \unsure{ds|dans} les \uline{fonctions rationnelles} de valeur $1$ en $0$}
qui \inserted{(composé avec l'\deleted{\unsure{involution|antiautomorphisme}}\inserted{aut} ``additif'' $\varphi(t)\longmapsto\frac{1}{\varphi(t)}$ dans $1+k[[t]]^+$)} est un hom de $\lambda$-anneaux,

Lorsque $k$ est \unsure{un|\deleted{un}} corps, l'hom précédent est injectif. En effet, \deleted{si} \deleted{$k$ est alg. clos,} la connaissance de la série formelle permet de reconstituer \inserted{la fonction rationnelle} \inserted{\unsure{donc}} [$k[t]$ étant factoriel.] l'élément $\sum_{\substack{P\ \mathrm{pol.}\\\text{unitaire}}}c_P(V_P,f_P)$ \deleted{dont} par
\[
L=\prod_Q Q^{c_P}
\]
[$Q$ et $P$ associés de la façon précisée plus haut].

Notons qu'on a un homomorphisme naturel \deleted{\unsure{le groupe log}}
\[
\begin{aligned}
1+k[[\unsure{1|t}]]^+&\longrightarrow k[[t]]\\
\varphi(t)&\longmapsto\frac{\varphi'(t)}{\varphi(t)}\deleted{\unsure{=}}
\end{aligned}
\]
qui correspond à
\[
\Lambda(V,f)(t)=\frac{L(V,f)'(t)}{L(V,f)(t)}
=\bigl[\log L(V,f)\bigr]'=\frac{1}{t}\sum_{1}^{\infty}(\operatorname{Tr}f^n)t^n.
\]
\deleted{\unsure{En car. 0, on peut|En car. 0, on}} Si $k$ de car. $0$, \inserted{i.e. $k\supset\mathbb{Q}$} on peut inverser cet homomorphisme par $\exp\bigl(\int_0^{\deleted{\unsure{t}}\ldots}\Lambda(V,f)(t)\bigr)$. Donc \inserted{si $k$ est} \inserted{un corps de car. $0$,}
\[
\Lambda:R(\mathbb{N},k)\longrightarrow k[[t]]
\]
est alors injectif \deleted{($k$ corps}
% archive sheet number: 90
